\summarychapter{36}{The Restoration of Fortresses}{Catholicism in the Age of Progress}

\begin{inbrief}
The fall of the Bastille opened the Age of Progress (1789--1914), whose creed was human progress and
whose watchwords were liberty, equality and fraternity. The Roman Catholic Church, bound to the old
order, lost much in the storm and answered by trying to turn back the clock. The chapter follows the
French Revolution's attack on the church, the popes' war with liberalism, the loss of the Papal
States and the First Vatican Council's definition of papal infallibility.
\question{Why did Catholicism look with such fear upon the popular movements of the time?}
\end{inbrief}

On the eastern edge of Paris stood the Bastille, an old fortress used as a prison, which propaganda
had made a symbol of royal tyranny. On the morning of July 14, 1789, a crowd gathered before it.
Patriots cut the drawbridge chains, and Governor de Launay, with 110 soldiers, agreed to surrender
in return for safe withdrawal. When the inner gates opened the mob rushed in and murdered him. The
dungeons held seven victims of tyranny: five ordinary criminals and two madmen. The sordid facts were
soon turned into heroic legend.

\ssection{Birth of a New Age}
The Bastille stood for the Old Regime: absolute monarchs and a feudal society of church, aristocrats
and powerless commoners. The crowd stood for the nineteenth century and the rights of the common man.
Few doubted that history moves upward and that the race is growing better and happier. Christianity
made its way through the period under adverse conditions, and Rome, long tied to the old order, saw
many of its treasures swept away. The democratic gospel of the Revolution glorified man rather than
God, and Rome saw more clearly than most Protestants that the devil can be democratic when it suits
him. Ten thousand people telling a lie do not make it true, Shelley remarks. Votes and schooling do
not bring utopia, because the flaw in human nature lies deeper than any fault in institutions.
Alexis de Tocqueville, in \emph{Democracy in America}, warned that equality itself may produce a
``tyranny of the majority,'' rule by shifting public moods. Unhappily Rome, in resisting the gospel
of the common man, tried to build a medieval fortress on the road of progress, and the secular masses
passed by on the other side.

\emph{Liberty} meant individual freedom in politics and economics, though the word is used in
confusing ways, by Ronald Reagan and Ted Kennedy alike, who differed about what it looks like and how
to reach it. The early liberals spoke for the middle class. They wanted the vote, representative
government, and freedom to build factories and gather wealth without interference (\emph{laissez
faire}). \emph{Equality} meant the rights of men whatever their birth or wealth. Peasants and urban
workers sought equality with the middle class through socialism, to be reached by evolution within
democracy or by revolution on the Marxist pattern. \emph{Fraternity} meant the sense of brotherhood
that became nationalism, the ambition of a people to be masters of their own territory and destiny,
which swept Europe in the nineteenth century and Asia and Africa in the twentieth.

On the eve of the Revolution the church basked in the old order, which it had sanctified for a
thousand years, blessing kings and nobles and giving little thought to peasants or the rising middle
class. Of France's twenty-five million people only two hundred thousand, nobles and clergy, were
privileged, yet they controlled nearly half the land and the best offices. The peasants, 80 percent
of the population, staggered under taxes to church and state. The middle class had wealth without
responsibility, intelligence without authority and ability without recognition.

\ssection{The Revolutionary Fever}
The Enlightenment had created a ferment, mostly of words. From the 1760s political unrest spread
through states small and large, with the same demands everywhere: participation, the vote, freedom of
expression. The American Revolution inspired Europe's radicals. The colonists seemed true men of the
Enlightenment, and their success proved its ideas in the laboratory of experience. In France
bankruptcy loomed. The government borrowed heavily and falsified its accounts, churchmen lived high,
and the wine harvests failed. In 1789 Louis XVI summoned the Estates-General, the assembly of clergy,
nobility and commons, which had not met for 175 years. The upper estates wanted to vote by estate,
which would keep control with them; the Third Estate, now half the membership, wanted to vote by
head. When the king refused, the commoners broke away and formed the National Assembly. After July 14
the crown could no longer keep order. By the end of August most feudal privileges were gone and the
``Declaration of the Rights of Man and of the Citizen'' was law. It declared the natural rights of
``liberty, property, security, and resistance to oppression'' sacred, guaranteed free expression,
forbade arbitrary arrest, and proclaimed France a sovereign nation that belongs to its people, not to
its monarchs. Within ten years France had formed a republic, executed its king, passed through a
revolutionary regime and a time of confusion, and ended with the coup that brought General Napoleon
Bonaparte to power, all the while at war with the rest of Europe.

The church was a special object of rage. In the early 1790s the National Assembly reformed it on
Enlightenment lines, giving pastors a decent income and redrawing dioceses, which was to the good.
But when it abolished papal control and imposed a loyalty oath, it split the church. In almost every
town the constitutional clergy who swore faced the non-constitutional who refused, and thirty to
forty thousand priests were driven into exile or hiding. The Revolution then took on a religion of
its own. A new calendar removed every trace of Christianity. Parish churches became Temples of
Reason, and in Notre Dame an actress was enthroned on the high altar as the Goddess of Reason, while
in the provinces girls dressed as Reason, Liberty or Nature led processions. By 1794 the parody was
spent, and a decree early the next year guaranteed the free exercise of religion; Catholics returned
to the altar. But Rome never forgot that liberty had meant the worship of the Goddess of Reason.
Napoleon's Concordat of 1801 restored Catholicism to a special place as ``the religion of the great
majority of Frenchmen,'' but its power was gone for good. Europe would not return to the alliance of
throne and altar, and Rome never acquired a taste for liberalism.

\ssection{Catholicism versus Liberalism}
Jesus and the apostles said little about personal independence or the right to one's opinions.
Through the Middle Ages and the Reformation Augustine's axiom ruled: liberty comes by grace, not
grace by liberty. To be truly free a person must be in a state of salvation, so Christians had felt
little enthusiasm for freedom in the merely political sense. In the nineteenth century the idea was
in the air that everyone should be as free as possible.

\begin{sidebox}{Liberty in an Ancient Sense}
In many ancient settings liberty meant more than an individual's power to choose. It was tied to
nature. Human beings have a nature or essence, and liberty meant being able to live up to it. A man
crawling in slop would win no praise on the ground that he had chosen to act like a pig. Freedom was
noble because it allowed a person to reach his potential, and barriers to that should be removed.
Today people often understand liberty apart from nature, as the power to choose what they want to be.
In the same way the ancients tied happiness to being or becoming good, not just to feeling good, and
Aristotle would have found it odd to be content that someone feels good without being good. The
American founders assumed this older idea when they spoke of life, liberty and happiness. Present
notions of liberty are often poor beside it.
\end{sidebox}

John Stuart Mill defined the possible: ``the liberty of each, limited by the like liberty of all.''
Everyone may hold, spread and act on his opinions so far as others may do the same. In practice this
meant constitutional government with civil liberties for all, including freedom of worship. The popes
despised the thought. After Napoleon's fall in 1815 absolutist kings tried to recover their ground
and met liberal resistance everywhere. Risings in Spain and Italy were put down, but the restored
French kings were overthrown in 1830, and in 1848 liberal revolution briefly triumphed in almost
every capital. The popes of these years, Leo XII, Pius VIII and Gregory XVI, were not bad men, but
they refused to join their century. They defended the past and never learned how to fight the new
world or how to convert it. Liberalism for its part refused the church's help and denied it any right
to speak on the morality of public life: politics is independent of Christian ethics, and Catholics
are private citizens and nothing more.

The plainest sign of the papacy's tie to the past was the Papal States, where the pope was an earthly
sovereign. Italy had for centuries been only a geographical expression, seven states besides the
papal territory that ran northeast from Rome across the peninsula. In mid-century the
\emph{risorgimento} (``rebirth''), arising in Sardinia, aimed to expel foreign powers and unite the
peninsula, and it could not tolerate a medieval, absolutist state in the middle. After 1849 the Papal
States were so hated that only French bayonets could hold them. Liberals at first welcomed Pius IX
(1846--1878), a warm and kindly man who on March 14, 1848, gave his states a constitution; some
dreamed of an Italian federation under the pope. Then revolutionaries assassinated his prime
minister, Count Pellegrino Rossi, revolution broke out in Rome, and Pius fled. Restored by French
troops, he returned to absolutist rule. The national movement under Victor Emmanuel II of Sardinia
became an avalanche. In 1859--1860 most of the Papal States fell, and in March 1861 Victor Emmanuel
was proclaimed king of Italy. When the Franco-Prussian War drew off the French garrison, the
nationalists bombarded Rome, which surrendered on September 20, 1870. After more than a thousand
years the Papal States were at an end.

Pius withdrew into the Vatican. The king moved to Rome in June 1871, ignoring protests and
excommunications. The government offered an annual subsidy and the free exercise of all spiritual
functions; Pius angrily refused and called himself the ``prisoner of the Vatican.'' He forbade
Italian Catholics to take part in elections, which only left the field to the radicals and made the
government more anticlerical. This ``Roman Question'' was settled only in February 1929, when
Mussolini concluded the Lateran Treaty: the pope gave up all claim to the former states and received
full sovereignty in the small Vatican State.

\ssection{The Infallibility of the Pope}
The year 1870 saw both the end of the pope's earthly rule and the declaration of his supreme
authority and infallibility. The First Vatican Council was the climax of ultramontanism, devotion to
Rome ``across the mountains.'' After the upheavals of revolution and empire some French Catholics
exalted the papacy as the only source of civil order and public morality. Only a clergy independent
of the state and led by an infallible pope, undisputed master of the church, could protect spiritual
freedom from political tyranny. The church must be a monarchy, and infallibility is to the pope what
sovereignty is to a king. By mid-century this had a large following, which the pope encouraged. A
Jesuit journal explained that when the pope meditates, God is thinking in him; hymns were addressed
to Pius IX, and some called him ``the vice-God of humanity.''

On December 8, 1854, Pius declared as dogma the traditional belief in the immaculate conception:
Mary was from the first moment of her conception preserved by special grace from the stain of
original sin. The belief was not new. The manner was: fifty-four cardinals and 140 bishops attended,
but the decision was the pope's alone, a definition \emph{ex cathedra}, ``from the chair.'' Could the
pope define dogma without a council? Ten years later, in December 1864, as the nationalists closed in
on his states, Pius sent the bishops an encyclical with a \emph{Syllabus of Errors} listing eighty
evils of modern society. He declared war on socialism, rationalism, freedom of the press, freedom of
religion, public schools, Bible societies and the separation of church and state, and denied that the
Roman pontiff ``can and ought to reconcile himself'' to ``progress, liberalism, and modern
civilization.'' There would be no peace talks. The church must close ranks around its leader.

To strengthen his hand Pius prepared a general council. He appointed a preparatory body of cardinals
in 1865, announced his plan to more than five hundred bishops in 1867, at the eighteen-hundredth
anniversary of the martyrdom of Peter and Paul, and opened the council in Rome on December 8, 1869.
Catholics did not doubt that Peter's successor has special teaching authority. The question was how
far it extends, whether it can be used apart from councils and the college of bishops, and on what
conditions. The opposition to a definition was led by Bishop Hefele of Rottenburg, the learned
historian of the councils, and Bishop Strossmayer of Djakovar, supported by many others, including
most of the Germans. In the first ballot, on July 13, 1870, 451 voted for the definition, 88 against
and 62 with reservations; many opponents thought only that the time was wrong. To avoid scandal
fifty-five bishops left Rome, with the pope's consent, before the final vote on July 18: 533 in
favor and two against.

The council thus asserted two things. First, primacy: as successor of Peter, vicar of Christ and
supreme head of the church, the pope has full and direct authority over the whole church and over
each bishop, in faith and morals and also in discipline and administration. Second, infallibility:
when the pope in his official capacity, \emph{ex cathedra}, makes a final decision for the whole
church in a matter of faith and morals, that decision is of itself infallible and unchangeable and
needs no prior consent of the church. The Franco-Prussian War, which broke out the next day, sent
many fathers home, and the occupation of Rome ended any hope of resuming. The vital work was done.
The council had, so to speak, restored the Bastille. Surrounded by liberalism, socialism and
nationalism, Rome chose safety behind the walls of an exalted, infallible papacy, a strategy that
shaped Catholic life for generations. But fortresses, Shelley concludes, grow stuffy. They allow no
enlargement of thought, and those inside come to imagine that the only world that matters lies
within the walls.

\begin{keyfacts}
\keypeople{Louis XVI; Napoleon Bonaparte; Alexis de Tocqueville; John Stuart Mill; Pius IX; Count
Pellegrino Rossi; Victor Emmanuel II; Bishop Hefele; Bishop Strossmayer; Mussolini.}
\keydates{%
  \kd{1789}{Fall of the Bastille (July 14); Declaration of the Rights of Man}
  \kd{1801}{Napoleon's Concordat with the pope}
  \kd{1848}{Liberal revolutions across Europe; Pius IX flees Rome}
  \kd{1854}{Pius IX defines the immaculate conception of Mary}
  \kd{1861}{Victor Emmanuel II proclaimed king of Italy}
  \kd{1864}{The \emph{Syllabus of Errors}}
  \kd{1869--1870}{First Vatican Council: papal primacy and infallibility}
  \kd{1870}{Rome falls to the Italian nationalists (September 20); end of the Papal States}
  \kd{1929}{Lateran Treaty creates the Vatican State}}
\keyterms{Age of Progress; Old Regime; liberty, equality, fraternity; liberalism; \emph{laissez
faire}; socialism; nationalism; cult of Reason; Concordat; \emph{risorgimento}; Papal States;
``prisoner of the Vatican''; ultramontanism; \emph{ex cathedra}; infallibility; \emph{Syllabus of
Errors}.}
\end{keyfacts}
